% A synthetic showyourwork paper: syw-site's example and end-to-end test. Not real science.
\documentclass[twocolumn]{aastex631}
\usepackage{showyourwork}
\let\tablenum\relax % AASTeX 6.3.1 defines \tablenum, as siunitx does
\usepackage{siunitx}
\DeclareSIUnit\kpc{kpc}

\begin{document}

\title{A Helix Seen From Every Side: an Example Paper for \texttt{syw-site}}

\author[0000-0002-1825-009X]{Ada Example}
\affiliation{Department of Examples, Example University}
\author{Bo Sample}
\affiliation{Department of Examples, Example University}
\affiliation{Institute for Placeholders}

\begin{abstract}
We draw a helix of $\variable{output/turns.tex}$ turns and measure its mean projected radius,
\variable{output/radius.tex}, from above. Fully 100\% of the numbers here are made up.
\end{abstract}

\section{The helix}
The helix has \variable{output/turns.tex} turns and a pitch of \variable{output/pitch.tex}
(Figure~\ref{fig:helix}); seen side on it is a wave (Figure~\ref{fig:sides}).

\begin{figure}[ht!]
    \script{helix.py}
    \begin{centering}
        \includegraphics[width=\linewidth]{figures/helix.pdf}
        \caption{The helix from above, $r = 1$.}
        \label{fig:helix}
    \end{centering}
\end{figure}

\begin{figure*}[ht!]
    \script{sides.py}
    \begin{centering}
        \includegraphics[width=0.48\linewidth]{figures/side_x.pdf}
        \includegraphics[width=0.48\linewidth]{figures/side_y.pdf}
        \caption{The helix seen along $x$ (left) and $y$ (right), as in Figure~\ref{fig:helix}.}
        \label{fig:sides}
    \end{centering}
\end{figure*}

\end{document}
